\section{Composable E-Graph Extensions}
\label{sec:extensions}

% Sources: ../extended-egglog-paper/abstract.tex and intro.tex;
% submission status supplied by Oliver.
\sourcenote{Edited from \emph{Extended E-Graphs via Compilation}, abstract
  and introduction; new connection through the running example.}
The preceding sections address the size of proofs and the choice of an executable program.
Producing proofs for a realistic optimizer also requires support for analyses and variable binding.
In the running summation example, an analysis establishes that a factor is independent
  of the bound variable, and the proof uses that fact to justify moving the factor.
Renaming the sum's variable also changes the analysis fact that the proof must justify.
These interactions make independently implemented extensions difficult to combine.

% Source: ../extended-egglog-paper/intro.tex, architecture and contributions.
\sourcenote{Edited from \emph{Extended E-Graphs via Compilation}, introduction;
  submission status added.}
Our work, \emph{Extended E-Graphs via Compilation}, is currently under submission~\cite{extended-egglog}.
It implements slotted e-graphs, proof production, and disequality reasoning
  as compiler passes over \texttt{egglog} programs.
Disequality reasoning derives explicit facts that terms have different values.
Each pass translates an extension into rules in a simpler language,
  making interactions explicit at the boundary between passes.
This section focuses on composing binding, analyses, and proofs;
  the complete translations and correctness arguments appear in the paper.

\subsection{Background: Variable Binding}
\label{sec:background-binding}

% Source: ../extended-egglog-paper/slotted.tex, opening and Semantics;
% examples adapted to continue the report's running summation example.
\sourcenote{Edited from \emph{Extended E-Graphs via Compilation},
  binding definitions; examples adapted.}
Binders such as summations and lambda abstractions introduce variables with local scope.
In $\sum_{i=0}^{10} i\,j$, the variable $i$ is bound by the summation,
  while $j$ is free.
Consistently renaming $i$ to $k$ yields $\sum_{k=0}^{10} k\,j$,
  an \emph{alpha-equivalent} expression with the same meaning.
Substitution respects these scopes so that a free variable does not become accidentally bound.

% Primary source: Schneider et al., PLDI 2025, abstract and Sections 2--3.
% https://steuwer.info/files/publications/2025/PLDI-Slotted-E-Graphs.pdf
\sourcenote{New summary of prior work: \emph{Slotted E-Graphs}, \S\S2--3.}
Slotted e-graphs share representations of terms that differ only
  in the names of bound or free variables~\cite{slotted-egraphs}.
Their e-classes use \emph{slots} as placeholders for free variables.
A reference to an e-class carries a renaming that connects its slots
  to variables in the surrounding term.
For example, the terms $i\,j$ and $k\,j$ can share a representation with two free slots,
  mapped to $i$ and $j$ in one reference and to $k$ and $j$ in the other.
These mappings preserve the distinction between the free variables.

\subsection{Extensions via Compilation}
\label{sec:extension-architecture}

% Sources: ../extended-egglog-paper/intro.tex and overview.tex, pass architecture.
\sourcenote{Edited from \emph{Extended E-Graphs via Compilation},
  introduction and overview.}
The compilation pipeline progressively removes features from the language
  interpreted by the engine (Figure~\ref{fig:extension-pipeline}).
The slotted pass replaces built-in variable renaming with rules over explicit renaming maps.
The congruence pass replaces built-in equality maintenance with rules that record
  and propagate equalities.
The proof pass then adds rules that log how facts were derived.
The resulting program runs in an equality-free core:
  it stores and queries terms, but all reasoning about equality takes place in generated rules.
These rules reuse \texttt{egglog}'s query optimizations and parallel execution.

% Source: ../extended-egglog-paper/drawnfigures/architecture.tex, right-hand pipeline.
\begin{figure}[htbp]
  \sourcenote{New simplified diagram of the pipeline in
    \emph{Extended E-Graphs via Compilation}; new caption.}
  \centering
  % Simplified from ../extended-egglog-paper/drawnfigures/architecture.tex.
% New layout and labels; the pass order follows the paper.
\begin{tikzpicture}[
    font=\small,
    stage/.style={draw, rounded corners=2pt, align=center,
      minimum width=3.15cm, minimum height=1.1cm, inner sep=4pt},
    flow/.style={-{Latex[length=2.2mm]}, thick}]
  \node[stage, fill=orange!12] (slotted) at (0,0)
    {Slotted pass\\\footnotesize Renaming rules};
  \node[stage, fill=blue!8] (congruence) at (3.85,0)
    {Congruence pass\\\footnotesize Equality rules};
  \node[stage, fill=green!10] (proof) at (7.7,0)
    {Proof pass\\\footnotesize Logging rules};
  \node[stage, fill=black!4] (core) at (11.55,0)
    {Core engine\\\footnotesize No built-in equality};
  \draw[flow] (slotted) -- (congruence);
  \draw[flow] (congruence) -- (proof);
  \draw[flow] (proof) -- (core);
\end{tikzpicture}

  \caption{Compilation of an extended \texttt{egglog} program.
    Each pass expresses its feature using constructs supported by the next language.
    The proof pass logs both user-rule applications and generated maintenance steps;
    proof lifting translates those steps back into a certificate over the user's rules.
    An optional disequality pass precedes the slotted pass.}
  \label{fig:extension-pipeline}
\end{figure}

% Sources: ../extended-egglog-paper/overview.tex, schedule transformation;
% congruence.tex, Run translation and correctness.
\sourcenote{Edited from \emph{Extended E-Graphs via Compilation},
  overview and congruence pass.}
Generated rules must run before user rules rely on their results.
The congruence pass saturates the rules that maintain equality after each round of user rules,
  so the next round sees the equalities that built-in congruence closure would derive.
Our operational semantics specifies each scheduled step,
  allowing pass correctness to account for executions that stop before full saturation.

\subsection{Proofs for Equalities and Analysis Results}
\label{sec:extension-proofs}

% Source: ../extended-egglog-paper/proofs.tex, opening and Proofs;
% contrast with Section 3 added for the report.
\sourcenote{Edited from \emph{Extended E-Graphs via Compilation}, proof pass;
  new connection to Section~\ref{sec:proofs}.}
Section~\ref{sec:proofs} takes the equalities supplied to congruence closure as given.
For \texttt{egglog}, a certificate must also explain why a rule could supply each equality.
In the running example, moving $5$ inside the outer sum requires
  $\mathtt{NotFree}(\texttt{"i"},5)$.
The proof must derive this fact from the analysis rule,
  then use it as a premise of distribution.

% Sources: ../extended-egglog-paper/background.tex, initialization, analysis,
% and distribution rules; proofs.tex, existence claims and Rule proof steps.
\begin{samepage}
\sourcenote{New proof presentation of the running example from
  \emph{Extended E-Graphs via Compilation}; the inference steps follow its rules.}
Let $B=\sum_{j=0}^{10}i\,j$ and $S=\sum_{i=0}^{10}B$.
Write $E(t)$ for evidence that the term $t$ has been built; analysis facts are terms too.
The first distribution step has the following proof:
\[
  \dfrac{
    \dfrac{}{E(5\cdot S)}\;\text{insert}
    \qquad
    \dfrac{
      \dfrac{}{E(S)}\;\text{insert}
      \qquad
      \dfrac{}{E(5)}\;\text{insert}
    }{E(\mathtt{NotFree}(\texttt{"i"},5))}\;\text{analysis}
  }{5\cdot S=\displaystyle\sum_{i=0}^{10}5\cdot B}\;\text{distribution}.
\]
\end{samepage}
All three leaves use the same initialization rule, which inserts the input and its subterms.
The analysis branch uses the rule that numbers have no free variables.
The proof of distribution into the inner sum has the same structure.

% Source: ../extended-egglog-paper/proofs.tex, opening and Proofs.
\sourcenote{Edited from \emph{Extended E-Graphs via Compilation}, proof pass.}
The certificate format represents an existence claim $E(t)$ as $t=t$.
Although this equality is mathematically reflexive,
  it also asserts that the term $t$ has been derived in the \texttt{egglog} program.
An unconditional reflexivity step would allow a certificate to assert
  $\mathtt{NotFree}(\texttt{"i"},5)$ without evidence.
Our format therefore has no reflexivity step:
  existence claims must be justified by rule applications,
  just as other equalities must be justified by rules or equality reasoning.

% Source: ../extended-egglog-paper/proofs.tex, opening pipeline and provenance.
\sourcenote{Edited from \emph{Extended E-Graphs via Compilation}, proof pass;
  explanation of the pass order retained.}
To produce these certificates, the proof pass needs derivations for both equalities and analysis facts.
The congruence pass turns built-in equality maintenance into rule applications.
After this translation, every derived fact has the same kind of explanation:
  a rule firing and the facts that its query matched.
The proof pass logs these firings and reconstructs a proof by recursively explaining their premises.

% Sources: ../extended-egglog-paper/congruence.tex, Congruence Translation;
% proofs.tex, Lifting Proofs.
\sourcenote{Edited from \emph{Extended E-Graphs via Compilation},
  congruence pass and proof pass; path-compression rule unchanged apart from notation and layout.}
For example, the congruence pass represents links in the union-find structure,
  which tracks e-class membership,
  as terms $\mathtt{UF}(a,b)$ recording that $a$ and $b$ are equal.
One generated rule shortens a chain of links, an operation called \emph{path compression}:
\begin{eggprogram}
  &\mathtt{UF}(a,b)\land\mathtt{UF}(b,c)\land b\ne c&&\\*
  &\quad\Rightarrow\kw{subsume}(\mathtt{UF}(a,b)),\ \mathtt{UF}(a,c).&&
\end{eggprogram}
Here $\kw{subsume}$ removes the old link from future queries;
  the new link connects $a$ directly to $c$.

% Sources: ../extended-egglog-paper/proofs.tex, provenance and lifting;
% congruence.tex, path-compression rule above (unchanged apart from notation/layout).
\sourcenote{Edited from \emph{Extended E-Graphs via Compilation}, proof pass.}
The reconstructed proof describes the compiled program, including its generated rules.
To return a certificate for the user's original program,
  \emph{lifting} translates generated steps back into reasoning over the original rules.
For example, the path-compression step becomes transitivity of equality.

\subsection{Slotted E-Graphs via Compilation}
\label{sec:slotted-compilation}

% Source: ../extended-egglog-paper/slotted.tex, opening and Semantics.
\sourcenote{Edited from \emph{Extended E-Graphs via Compilation}, slotted pass.}
The slotted pass expresses variable renaming as rules,
  so the later proof pass can record its derivations.
Binding annotations specify which arguments introduce variables and where they are bound.
For summations, $\operatorname{bind}(\SumZeroTen)=\{(1,2)\}$ declares that the first argument
  names a variable bound within the second.
We write object-language variables as slots $\$i$ and $\$j$,
  distinguishing them from pattern variables such as $x$.
A renaming maps distinct slots to distinct slots and avoids capturing free variables.
Renaming both sides of an established equality yields another valid equality;
  changing only bound names identifies alpha-equivalent terms.
The generated rules operate over terms and explicit renaming maps.

% Source: ../extended-egglog-paper/slotted.tex, Slotted Maintenance Implementation.
\sourcenote{Edited from \emph{Extended E-Graphs via Compilation},
  slotted maintenance; argument pairs grouped for readability.}
The stored representation shares terms across renamings by pairing each child
  with a map from its slot names to the names used by the parent.
We write each child as a pair $(\text{renaming},\text{term})$.
All variable leaves can then use one stored term, $\$0$:
\[
  \mathtt{Mul}(\$i,\$j)
  \quad\leadsto\quad
  \mathtt{Mul}\bigl(([\$0\mapsto\$i],\$0),\,
                     ([\$0\mapsto\$j],\$0)\bigr).
\]
The maps keep $\$i$ and $\$j$ distinct despite sharing the leaf.
Generated rules discover terms that are equal up to renaming
  and choose a common stored representation for them.
Other generated rules maintain the free slots of each e-class,
  including when an equality makes a variable unnecessary.

% Source: ../extended-egglog-paper/slotted.tex, E-Matching Implementation.
\sourcenote{Edited from \emph{Extended E-Graphs via Compilation},
  slotted matching and proof lifting.}
Matching a source query against this representation must recover the relevant renamings.
The compiled query finds stored terms and explores compatible assignments of slot names.
Distinct slots within a term must remain distinct,
  while occurrences that the query equates must receive consistent names.
These constraints also apply across \emph{multipatterns},
  queries that match several patterns with shared variables.
Analyses use the same representation, so their facts participate in renaming
  and free-variable checks.
The later proof pass records the generated derivations;
  slotted proof lifting replaces maintenance steps with equality and renaming steps
  checked against the original rules and binding annotations.
Thus the combination supports binding, analyses, multipatterns, and proofs together.

\subsection{Correctness Guarantees}
\label{sec:extension-correctness}

% Sources: ../extended-egglog-paper/congruence.tex, Correctness;
% proofs.tex, Adapting Datalog Provenance; appendix.tex, Proof Checker.
\sourcenote{Edited from \emph{Extended E-Graphs via Compilation},
  congruence correctness, provenance correctness, and checker appendix.}
We formalize the core semantics, congruence pass, and proof-recording transformation in Lean.
For the formalized source fragment, these passes preserve
  whether a program succeeds, fails a check, or does not terminate.
We also prove that successfully lifted certificates are accepted by the source checker.
The checker's soundness theorem establishes that an accepted equality holds
  in an e-graph containing its endpoints and closed under the program's rules:
  rerunning any rule adds no new information.
This closure condition is needed because a certificate establishes a consequence
  of the rules independently of the user's schedule.

% Sources: ../extended-egglog-paper/proofs.tex, formalization scope;
% slotted.tex, final paragraph; appendix.tex, source-checker fragment.
\sourcenote{Edited from \emph{Extended E-Graphs via Compilation},
  proof pass and slotted pass; scope condensed.}
These theorems cover a restricted language, not the entire implementation.
The source checker excludes term removal and tests for whether terms currently occupy different e-classes.
It also excludes $\min$, which generated rules use to choose a common representative for equal terms.
They do not guarantee that proof reconstruction or lifting always succeeds.
Slotted compilation and its source checker remain outside the Lean formalization.

\subsection{Evaluation and Applications}
\label{sec:extension-evaluation}

% Source: ../extended-egglog-paper/eval.tex, research questions and application studies.
\sourcenote{Edited from \emph{Extended E-Graphs via Compilation}, evaluation questions.}
The evaluation tests whether compiled extensions preserve existing capabilities,
  how much proof production costs in existing optimizers,
  and whether combining extensions enables useful optimizations.

% Sources: ../extended-egglog-paper/eval.tex, RQ1--RQ2;
% eval-presentation.numbers.tex and figures/eval/headline-speedups.numbers.tex.
\sourcenote{Edited from \emph{Extended E-Graphs via Compilation}, evaluation.}
The generated rules reproduce the checked results of prior implementations.
The Math workload from the original \texttt{egglog} evaluation~\cite{egglog}
  also lets us compare the cost of producing and checking proofs:
  execution with proof extraction and checking takes 29.07 seconds,
  compared with 38.58 seconds for \texttt{egg} with proofs.
Proof simplification is disabled in both; our time includes compilation by the passes.

% Sources: ../extended-egglog-paper/eval.tex, Adding Proofs to Existing egglog Programs;
% eval-presentation.numbers.tex and figures/eval/corpus-proof-stages-cdf.numbers.tex.
\sourcenote{Edited from \emph{Extended E-Graphs via Compilation},
  application evaluation and proof-cost appendix; percentile measurements rounded.}
To test proof production in existing optimizers, we ran
  486 \texttt{egglog} programs from six projects, including \texttt{eggcc}.
Across these programs,
  we produced and checked a proof for every requested extraction or equivalence check.
Relative to original \texttt{egglog} execution, execution with proof recording takes
  less than $3.66\times$ as long for half the programs and less than $4.71\times$ for 95\%.
These ratios exclude proof reconstruction and checking.
Including reconstruction, simplification, and checking raises the median to $4.64\times$
  and the 95th percentile to $71.47\times$.
Inexpensive recording therefore does not by itself make validation inexpensive.

% Sources: ../extended-egglog-paper/eval.tex, LMFAO Query Optimization;
% appendix.tex, LMFAO Workload and Measurements.
\sourcenote{Edited from \emph{Extended E-Graphs via Compilation},
  LMFAO evaluation and measurement appendix.}
We test composition by implementing the LMFAO query optimizer~\cite{lmfao}
  with the combined passes.
Its rewrites move independent factors out of sums, interchange loops, and fuse reductions,
  requiring binding, analyses, and multipatterns together.
On 140 covariance queries over the released sample data, the slotted implementation produces
  checked proofs and queries with performance comparable to those chosen by the original optimizer,
  at higher planning and compilation cost.
An alternative encoding using numerical indices for bound variables
  reaches a 20\,GiB memory limit without producing a plan.
This case demonstrates both the benefit of composition and the cost of broader search.

% New synthesis from Sections 3 and 5 and Oliver's proposed compiler design.
\sourcenote{New: connection to proof minimization and the proposed compiler.}
Proof production for analysis facts supplies a component for the proposed compiler.
The measured checking costs motivate evaluating proof minimization
  on certificates containing analysis and renaming steps.
A further obligation is semantic: deriving a fact from \texttt{egglog} rules
  does not establish that those rules correctly describe program behavior.
Section~\ref{sec:proposed-research} proposes the remaining connection
  between accepted derivations and language semantics.

\subsection{Related Work}
\label{sec:related-extensions}

% Source: ../extended-egglog-paper/related_work.tex, Implementing Extensions by Compilation.
\sourcenote{Edited from \emph{Extended E-Graphs via Compilation}, related work.}
The nanopass approach offers an architectural precedent for combining language features:
  each pass removes a small language feature and produces a simpler intermediate language~\cite{nanopass}.
Ballantyne and colleagues similarly expand language extensions into a core language
  optimized by a nanopass back end~\cite{nanopass2}.
Applying this organization to e-graphs requires accounting for features
  that affect a shared equality relation:
  a change to variable binding or congruence can change what a rule matches
  and what evidence its application requires.

% Source: ../extended-egglog-paper/related_work.tex, Implementing Extensions by Compilation.
\sourcenote{Edited from \emph{Extended E-Graphs via Compilation}, related work.}
Relational e-matching recasts matching rewrite patterns as database joins,
  which combine matching records that agree on shared variables~\cite{relational-ematching}.
The \texttt{egglog} query engine builds on this approach~\cite{egglog}.
By expressing extension logic as rules, our passes reuse this query engine
  and its optimizations.

% Source: ../extended-egglog-paper/related_work.tex, Egraph Semantics.
\sourcenote{Edited from \emph{Extended E-Graphs via Compilation}, related work.}
Suciu and colleagues characterize the result of running rewrite rules
  until no new terms or equalities can be added~\cite{semantic-foundations}.
This characterizes saturation independently of a particular implementation.
Our operational semantics describes scheduled executions,
  including runs that stop before saturation.

% Source: ../extended-egglog-paper/related_work.tex, Proof Production and Provenance.
\sourcenote{Edited from \emph{Extended E-Graphs via Compilation}, related work.}
Database \emph{provenance} records which input data and derivation steps contribute to a result.
Datalog and SQL systems obtain provenance through program or query transformations
  that log dependencies~\cite{datalog-debugging,exspan,provsql}.
Our proof pass applies this approach after congruence compilation
  and lifts the resulting derivations into proofs over the source rules.
